% Restrained A4 academic examination layout; non-official practice edition.
\usepackage{xeCJK}
\usepackage{fontspec}
\setmainfont{Noto Serif CJK SC}
\setsansfont{Noto Sans CJK SC}
\setmonofont{Noto Sans Mono CJK SC}
\setCJKmainfont{Noto Serif CJK SC}
\setCJKsansfont{Noto Sans CJK SC}
\setCJKmonofont{Noto Sans Mono CJK SC}
\usepackage{microtype}
\usepackage{fancyhdr}
\usepackage{lastpage}
\usepackage{needspace}
\usepackage{xcolor}
\usepackage{tabularx}
\usepackage{array}
\usepackage{booktabs}
\usepackage{longtable}
\usepackage{fvextra}
\usepackage{graphicx}
\usepackage{etoolbox}
\definecolor{examInk}{HTML}{20252B}
\definecolor{examGray}{HTML}{68727B}
\definecolor{examLine}{HTML}{BCC3C8}
\definecolor{examShade}{HTML}{F3F4F5}
\color{examInk}
\setlength{\parindent}{0pt}
\setlength{\parskip}{2.4pt plus .3pt minus .3pt}
\linespread{1.15}
\setlength{\emergencystretch}{2.6em}
\renewcommand{\arraystretch}{1.12}
\setlength{\tabcolsep}{4pt}
\setlength{\LTpre}{5pt}
\setlength{\LTpost}{6pt}
\setlength{\headheight}{16pt}
\pagestyle{fancy}
\fancyhf{}
\fancyhead[L]{\sffamily\footnotesize 计算机学科专业基础综合\quad\textcolor{examGray}{/}\quad 原创模拟卷 01}
\fancyhead[R]{\sffamily\footnotesize\textcolor{examGray}{\ifdefined\answeredition 解析版\else 学生卷\fi}}
\fancyfoot[L]{\sffamily\scriptsize\textcolor{examGray}{非官方模拟试题\,·\,仅供学习与测评}}
\fancyfoot[R]{\sffamily\scriptsize\textcolor{examGray}{\thepage\ / \pageref{LastPage}}}
\renewcommand{\headrulewidth}{0.30pt}
\renewcommand{\footrulewidth}{0pt}
\renewcommand{\headrule}{\hbox to\headwidth{\color{examLine}\leaders\hrule height \headrulewidth\hfill}}
\usepackage{titlesec}
\titleformat{\section}{\sffamily\bfseries\large}{}{0pt}{}
\titleformat{\subsection}{\sffamily\bfseries\normalsize}{}{0pt}{}
\titleformat{\subsubsection}{\sffamily\bfseries\normalsize}{}{0pt}{}
\titlespacing*{\section}{0pt}{11pt}{6pt}
\titlespacing*{\subsection}{0pt}{7pt}{4pt}
\titlespacing*{\subsubsection}{0pt}{6pt}{3pt}
\setlength{\fboxsep}{6pt}
\newcommand{\StudentCover}{%
 \thispagestyle{empty}%
 \noindent\begin{minipage}{\linewidth}\centering
 {\sffamily\fontsize{8.5}{12}\selectfont\color{examGray}\textbf{2026 · ORIGINAL PRACTICE EXAM}}\\[7pt]
 {\sffamily\bfseries\fontsize{19}{26}\selectfont 计算机学科专业基础综合}\\[3pt]
 {\sffamily\fontsize{12}{18}\selectfont 原创模拟试卷（一）}\\[13pt]
 {\color{examInk}\rule{\linewidth}{1.3pt}}\\[-6pt]
 {\color{examLine}\rule{\linewidth}{0.35pt}}\\[6pt]
 {\sffamily\fontsize{10}{15}\selectfont 考试时间：\textbf{180 分钟}\quad\;\;\;试卷满分：\textbf{150 分}\quad\;\;\;试卷共 \textbf{47} 题}
 \end{minipage}\par\vspace{10pt}
 \noindent\begin{tabularx}{\linewidth}{@{}lXlX@{}}
 \sffamily 姓名： &\makebox[\linewidth]{\dotfill}&\sffamily 考号： &\makebox[\linewidth]{\dotfill}
 \end{tabularx}\par\vspace{11pt}
 \noindent\begin{minipage}{\linewidth}
 \color{examInk}\sffamily\fontsize{9.1}{14}\selectfont
 \textbf{考生须知}\quad ① 全卷分单项选择题和综合应用题两部分；② 单项选择题每题仅有一个正确选项；③ 综合题应写出必要推导步骤，等价正确方法可得分。\par
 \vspace{3pt}\color{examGray}约定：除题目另有说明，H 表示十六进制，KiB = $2^{10}$ B，Mb/s = $10^{6}$ bit/s。
 \end{minipage}\par\vspace{5pt}}
\newcommand{\TeacherCover}{%
 \thispagestyle{empty}%
 \noindent\begin{minipage}{\linewidth}\centering
 {\sffamily\fontsize{8.5}{12}\selectfont\color{examGray}\textbf{SET 01 · TEACHER EDITION}}\\[8pt]
 {\sffamily\bfseries\fontsize{18}{25}\selectfont 参考答案与详细解析}\\[4pt]
 {\sffamily\fontsize{11}{16}\selectfont 计算机学科专业基础综合 · 原创模拟试卷（一）}\\[12pt]
 {\color{examInk}\rule{\linewidth}{1.3pt}}\\[-6pt]
 {\color{examLine}\rule{\linewidth}{0.35pt}}
 \end{minipage}\par\vspace{9pt}
 {\sffamily\fontsize{9.2}{14}\selectfont\color{examGray}本册与学生卷逐题对应。选择题给出判定依据，综合题提供主要推导及分步评分建议；阅卷时应认可等价正确解法。}\par\vspace{9pt}}
\newcommand{\ExamSection}[2]{%
 \par\Needspace{5\baselineskip}\vspace{10pt}
 \noindent\begin{tabularx}{\linewidth}{@{}>{\raggedright\arraybackslash}Xr@{}}
 {\sffamily\bfseries\fontsize{12}{16}\selectfont #1} & {\sffamily\color{examGray}\fontsize{9}{13}\selectfont #2}
 \end{tabularx}\par\vspace{-4pt}%
 {\color{examInk}\rule{\linewidth}{0.7pt}}\par\vspace{4pt}}
\newcommand{\TeacherSection}[1]{%
 \par\Needspace{5\baselineskip}\vspace{11pt}
 \noindent{\sffamily\bfseries\fontsize{11.5}{16}\selectfont #1}\par\vspace{2pt}%
 {\color{examLine}\rule{\linewidth}{0.45pt}}\par\vspace{2pt}}
\newcommand{\ExamQ}[1]{\par\Needspace{5\baselineskip}\vspace{6.5pt}\noindent{\sffamily\bfseries #1.}\hspace{.55em}}
\newcommand{\BigQ}[2]{%
 \par\Needspace{11\baselineskip}\vspace{10pt}\noindent
 \begingroup\setlength{\fboxsep}{6pt}
 \colorbox{examShade}{\parbox{\dimexpr\linewidth-2\fboxsep\relax}{%
 \sffamily\bfseries 第 #1 题\hfill\normalfont\sffamily\color{examGray}#2 分}}
 \endgroup\par\vspace{6pt}}
\newcommand{\Choice}[2]{\textsf{\bfseries #1.}\hspace{.4em}#2}
\newcommand{\ChoicesTwo}[4]{%
 \par\vspace{1pt}\noindent
 \begin{tabularx}{\linewidth}{@{}>{\raggedright\arraybackslash}X@{\hspace{7mm}}>{\raggedright\arraybackslash}X@{}}
 \Choice{A}{#1} & \Choice{B}{#2}\\[2pt]
 \Choice{C}{#3} & \Choice{D}{#4}
 \end{tabularx}\par\vspace{1pt}}
\newcommand{\ChoicesOne}[4]{%
 \par\vspace{1pt}\noindent
 \begin{tabularx}{\linewidth}{@{}>{\raggedright\arraybackslash}X@{}}
 \Choice{A}{#1}\\[2pt]
 \Choice{B}{#2}\\[2pt]
 \Choice{C}{#3}\\[2pt]
 \Choice{D}{#4}
 \end{tabularx}\par\vspace{1pt}}
\newcommand{\TeacherItem}[2]{%
 \par\Needspace{4\baselineskip}\vspace{6.5pt}\noindent
 {\sffamily\bfseries #1.}\hspace{.7em}
 {\sffamily\color{examGray}\fontsize{9}{12}\selectfont 答案}\,\textbf{#2}\par\vspace{-1pt}}
\newcommand{\TeacherBig}[2]{%
 \par\Needspace{8\baselineskip}\vspace{11pt}\noindent
 \begingroup\setlength{\fboxsep}{6pt}
 \colorbox{examShade}{\parbox{\dimexpr\linewidth-2\fboxsep\relax}{%
 \sffamily\bfseries 第 #1 题\hfill\normalfont\sffamily\color{examGray}评分总额 #2 分}}
 \endgroup\par\vspace{6pt}}
\newcommand{\ExamFigure}[3]{%
 \par\Needspace{7\baselineskip}\vspace{5pt}
 \begin{center}\includegraphics[width=#2\linewidth,height=0.27\textheight,keepaspectratio]{#1}\\[3pt]
 {\sffamily\fontsize{8.5}{11}\selectfont\color{examGray}#3}\end{center}\vspace{2pt}}
\newcommand{\AnswerGrid}[1]{\par\vspace{3pt}\noindent #1\par\vspace{6pt}}
\AtBeginDocument{\raggedbottom}
